\scanpage{25}
\mspage{12}
to examine the actions of $\pi_1(f),\pi(g)$ on the abelianized groups of the $\pi_1$, and even on their $\ell$-adic components. Taking then the generalized Jacobian of the type ``extension of an abelian variety by a torus'' of $X$, we know that

1°) The morphisms $f:X\longrightarrow G$ such that $f(a)=0$ factor uniquely through
\diagram{\begin{tikzcd}[ampersand replacement=\&]X \arrow[r,"{\mathrm{can}}"] \& J \arrow[r,"\varphi"] \& G\end{tikzcd}},
with $\varphi$ a homomorphism of algebraic groups

2°) \deleted{$\varphi$ is} Such a homomorphism $\varphi$ is known when one knows its action on the $H_1(\ ,\mathbb Z_\ell)$ Which amounts to knowing it on the points of order $\ell^\nu$, whatever $\nu$ may be -- now these are dense\ldots

3°) $H_{\unsure{1|n}}(X,\mathbb Z_\ell)\unsure{\simeq|\overset{\sim}{\longrightarrow}}H_1(J,\mathbb Z_\ell)$

From this, \deleted{the \unsure{case}} we conclude \inserted{(by 3°)} that $H_1(f)=H_1(g)$ implies (if $f=\varphi\circ\mathrm{can}$, $g=\psi\circ\mathrm{can}$) $H_{\unsure{1|n}}(\varphi)=H_1(\psi)$, hence by 2°) that $\varphi=\psi$, hence $f=g$, which was to be proved

\deleted{b)} Note that one cannot conclude from the hypothesis \inserted{of a)} that $f=g$, since for $X=Y=G$, the translations $\tau_a$ induce the identity on $\pi_1(G)$ (at least \unsure{in|\deleted{in}|\deleted{if}} characteristic $0$ or if $G$ is an abelian variety), without being the identity.

b) We are nevertheless going to prove the equality,
\margin{24}
